\clearpage\begingroup
\SourceLeaf{172}\HistoricalHead{167}{R}{GAUSS AN OLBERS.}
\begin{center}
[24.]\par\medskip
\textbf{\textsc{Gauss} an \textsc{Olbers}.\quad Göttingen, 18. Februar 1832.}\par\medskip
\rule{.78\linewidth}{.4pt}\par\smallskip
{\footnotesize [\textsc{Wilhelm Olbers}, Sein Leben und seine Werke II, 2, 1900, S. 584.]}\par\smallskip
\rule{.78\linewidth}{.4pt}
\end{center}
\noindent . . . . . . . . . . . . . . . . . . . . . .\par
\noindent\hspace*{\parindent}Vielleicht macht es Ihnen eine kleine Zerstreung, wenn ich Ihnen meine Grundidee zu einem Photometer, wie ich sie \textsc{Gerling} angegeben habe, anzeige. Denken Sie sich einen Spiegelsextanten mit der Modification, dass der kleine Spiegel gar nicht belegt ist, sondern bloss von seinen Glasflächen reflectirt, beide Spiegel aber \textls[120]{reichlich} so gross wie das Objectiv des Fernrohrs, der grosse auch so breit, wie es für die äussersten Fälle der Winkeldistanz zwischen zwei zu vergleichenden Sternen nöthig ist, damit \textls[120]{jeder} Punkt des Objectivs Licht bekomme. Man stellt den Sextanten so, als wollte man jene Distanz messen, so dass beide Bilder nahe beieinander erscheinen. Die ursprünglichen Lichtintensitäten der Sterne $A,B$ seien $a$ und $b$, die Intensität des Lichtes der Bilder $\alpha a,\,\beta b$, wo $\alpha$ und $\beta$ von der Öffnung des Objectivs, der unvollkommenen Durchsichtigkeit der Gläser abhängen, $\beta$ ausserdem auch noch von der Angulardistanz der Sterne. Jenseits des kleinen Spiegels ist aber noch eine Vorrichtung angebracht, vermöge der man das direct gesehene Licht auf einen beliebigen Bruch $=\mu$ reduciren kann, indem man statt des vollen Objectivs nur einen Sektor $=\mu\times360^{0}$ Licht verstattet. Diess $\mu$ bestimmt man so, dass beide Bilder gleich hell erscheinen; man hat also
\[
\mu\alpha a=\beta b.
\]
\noindent Jetzt macht man einen zweiten Versuch, indem man den Sextanten umkehrt und also den vorher direct gesehenen Stern reflectirt sieht. In diesem zweiten Versuch trete $\mu'$ an die Stelle von $\mu$. Man hat also
\[
\mu'\alpha b=\beta a.
\]
\noindent War ursprünglich $a=b$, so wird man nothwendig $\mu'=\mu$ finden und vice versa; sind aber $a,b$ ungleich, so hat man
\[
a:b=\beta:\mu\alpha=\mu'\alpha:\beta=\sqrt{\mu'}:\sqrt{\mu}.
\]
\noindent Zugleich wird immer $\mu\mu'=\frac{\beta\beta}{\alpha\alpha}$ eine bloss von der Angulardistanz abhängige
\par\clearpage\SourceLeaf{173}\HistoricalHead{168}{L}{PHYSIK.\quad BRIEFWECHSEL.}
\noindent Grösse sein, über die man aus vielen Versuchen das Gesetz ausfindig machen kann; nachdem dies gefunden ist, werden auch einseitige Messungen ein Resultat geben. Dies ist die eigentlich mathematische Grundidee; es versteht sich, dass an die Stelle eines grossen Sextanten mit kleinem Fernrohr hier ein grosses Fernrohr mit kleinem Sextanten (nur um die Sterne bequem zugleich ins Gesichtsfeld bringen zu können) treten muss. . . . . . . . . . . . .
\clearpage\endgroup
